\section{Data Ingestion and Feature Building}\label{sec:ingestion}

\subsection{Source data and cleaning}
The source is a Synthea-generated cohort exported as CSV. Eleven tables are
loaded --- patients, organizations, providers, encounters, conditions,
medications, procedures, allergies, immunizations, care plans and observations
--- and consolidated into a single hospital so that the data resembles one
institution rather than a national network. Cleaning addresses the
characteristic artifacts of synthetic generation: trailing digits are stripped
from names, future-dated events are clamped to the present, encounter durations
are derived and capped per encounter class (e.g., wellness visits to a few
hours, inpatient stays to a month), and missing values are imputed by mode for
categorical fields and by median for numerical ones. Observations, the largest
table, are inserted in chunks and aggregated through a materialized view. The
result is a normalized relational schema with foreign keys linking every
clinical event back to a patient and, where applicable, an encounter.

\subsection{Patient-level feature engineering}
The analytical stages operate not on raw events but on \emph{one feature vector
per patient}, obtained by aggregating across all clinical tables. Numerical
features include age, the count of active conditions, the number of medications
and procedures, the number of distinct providers seen, lifetime medication days,
total encounter duration, and ratios such as a preventive-to-reactive care
ratio. Heavily skewed counts are log-transformed. A composite
\emph{clinical-severity score} is derived by standardizing
(\,procedures, medications, active conditions\,) and projecting them onto their
first principal component, compressing three correlated severity dimensions into
one axis.

Categorical features encode clinically meaningful states: gender, pregnancy
history, comorbidity flags (heart disease, diabetes, hypertension, and others)
obtained by keyword matching over condition descriptions, a bucketed condition
burden, a BMI category, and a \emph{dominant care-plan} category ranked by
clinical severity (from general rehabilitation up to oncology and palliative
care). Continuous vitals such as BMI, blood pressure and glucose are bucketed
into risk tiers so that they enter the model as interpretable categories rather
than raw numbers. The complete vector is cached as a JSON column on the patient
record to avoid recomputation in later stages.
